%% ch-advanced.tex — Advanced Topics for StarForth Users
%% SOURCE: .claude/CLAUDE.md, scraps/architecture/
%% Included from user-guide/main.tex

%%━━━━━━━━━━━━━━━━━━━━━━━━━━━━━━━━━━━━━━━━━━━━━━━━━━━━━━━━━━━━━━━━━━━━━━━━━
\part{Advanced Topics}
%%━━━━━━━━━━━━━━━━━━━━━━━━━━━━━━━━━━━━━━━━━━━━━━━━━━━━━━━━━━━━━━━━━━━━━━━━━

%%────────────────────────────────────────────────────────────────────────────
\chapter{The StarForth VM Internals for Users}
\label{chap:vm-internals}
%%────────────────────────────────────────────────────────────────────────────

You do not need to understand the VM implementation to use StarForth
productively. But knowing a little about what happens under the hood explains
several things you will observe in practice: why some words are faster than
others, what the physics diagnostics are measuring, and how the adaptive
runtime relates to the code you write.

%% ─────────────────────────────────────────────────────────────────────────
\section{How the Interpreter Loop Works}
\label{sec:interp-loop}

When you type a line at the \texttt{ok>} prompt, StarForth:

\begin{enumerate}
  \item Tokenizes the input, splitting on whitespace.
  \item For each token, searches the dictionary (hot-words cache first, then
    linear scan from the head of the dictionary linked list).
  \item If found and in interpretation mode: executes the word immediately.
  \item If found and in compilation mode (inside \texttt{: ...}): compiles a
    reference to the word into the current definition (unless the word is
    \texttt{IMMEDIATE}, in which case it executes at compile time).
  \item If not found: attempts to parse the token as a number and push it.
  \item If not a number: reports an error and aborts the current line.
\end{enumerate}

Execution heat (\loopref{1}) is updated on every word execution. The rolling
window (\loopref{2}) records each execution for statistical analysis. Both
happen transparently.

%% ─────────────────────────────────────────────────────────────────────────
\section{Direct Threading}
\label{sec:direct-threading}

In the \texttt{fastest} build, StarForth uses direct threading: each
compiled word body is a sequence of code pointers, and the inner interpreter
jumps from one to the next using computed-goto. This eliminates an indirection
level compared to indirect threading and accounts for a significant portion of
the performance difference between the \texttt{debug} and \texttt{fastest}
builds.

From a user's perspective, direct threading is invisible --- words behave
identically. The only observable effect is speed.

%% ─────────────────────────────────────────────────────────────────────────
\section{The Hot-Words Cache}
\label{sec:hot-words-cache}

StarForth maintains a hot-words cache: the most frequently executed words are
sorted to the front of the dictionary linked list. Dictionary search always
starts at the head, so hot words are found faster.

When you define a new word, the cache may be invalidated and rebuilt. For
short interactive sessions this is imperceptible. For long-running programs
that define many words at startup, you may notice a brief pause as the cache
rebuilds --- this is \loopref{1} doing its work.

\texttt{WORD-ENTROPY} shows the current heat values. Words you use heavily
will show higher heat values; words you have not called recently will show
lower values (heat decays over time via \loopref{3}).

%% ─────────────────────────────────────────────────────────────────────────
\section{Execution Tokens}
\label{sec:execution-tokens}

An execution token (XT) is an opaque value that identifies a word. Tick
(\texttt{'}) produces an XT:

\begin{lstlisting}[language=Forth]
ok> ' DUP .
12345 ok>    \ some dictionary address
\end{lstlisting}

The number printed is implementation-dependent --- it is a VM address, not
meaningful on its own. What you can do with an XT is pass it to
\texttt{EXECUTE} to run the word later:

\begin{lstlisting}[language=Forth]
ok> ' DUP EXECUTE   \ same as just typing DUP
\end{lstlisting}

Or store it for a callback:

\begin{lstlisting}[language=Forth]
VARIABLE HOOK
' NOOP HOOK !         \ default: do nothing

: SET-HOOK  ( xt -- )  HOOK ! ;
: RUN-HOOK  ( -- )     HOOK @ EXECUTE ;
\end{lstlisting}

XTs are also the key for ACL lookups (Chapter~\ref{chap:acl}).

%%────────────────────────────────────────────────────────────────────────────
\chapter{Performance Tuning Your Code}
\label{chap:perf-tuning}
%%────────────────────────────────────────────────────────────────────────────

%% ─────────────────────────────────────────────────────────────────────────
\section{Measuring Before Optimizing}
\label{sec:measure-first}

Use \texttt{WORD-ENTROPY} to identify which words are hot before attempting
to optimize. Optimizing a cold word wastes effort.

The profiler gives more detail. Enable it at the shell level:

\begin{lstlisting}[language=bash]
./build/amd64/standard/starforth --profile 1 --profile-report < myprogram.4th
\end{lstlisting}

The report shows each word ranked by execution count. Typically 80--95\% of
all executions concentrate in the top ten words.

%% ─────────────────────────────────────────────────────────────────────────
\section{Build Profiles and Their Effect}
\label{sec:build-profiles}

The choice of build profile has a larger effect than any FORTH-level
optimization for most programs:

\begin{center}
\begin{tabular}{lll}
\toprule
Build & Relative speed & When to use \\
\midrule
\texttt{debug}   & 1$\times$ (baseline) & Debugging \\
\texttt{all}     & 3--4$\times$ & Development testing \\
\texttt{fast}    & 5--6$\times$ & Normal production \\
\texttt{fastest} & 6--7$\times$ & Benchmarks, deployment \\
\texttt{pgo}     & 7--8$\times$ & Maximum throughput \\
\bottomrule
\end{tabular}
\end{center}

For most interactive use, \texttt{all} or \texttt{fast} is fine. For
benchmarking or deployment, use \texttt{fastest} or \texttt{pgo}.

%% ─────────────────────────────────────────────────────────────────────────
\section{Physics-Aware Programming}
\label{sec:physics-aware}

The adaptive runtime favors frequently-used words. If you have a word that is
called millions of times in a tight loop, you can expect it to migrate to the
hot-words cache quickly. Conversely, setup words called only at startup will
cool off and be pushed toward the tail of the dictionary list.

For reproducible benchmarks:

\begin{lstlisting}[language=Forth]
: BENCH-RUN  ( iterations -- )
    PHYSICS-RESET-STATS    \ clear heat history
    PHYSICS-FREEZE         \ stop adaptive updates
    0 DO
        MY-WORD-UNDER-TEST
    LOOP
    PHYSICS-THAW ;
\end{lstlisting}

\texttt{PHYSICS-FREEZE} ensures the cache and heat values do not change during
measurement. \texttt{PHYSICS-THAW} restores normal operation afterward.

%% ─────────────────────────────────────────────────────────────────────────
\section{Stack Manipulation Cost}
\label{sec:stack-cost}

Stack operations (\texttt{DUP}, \texttt{SWAP}, \texttt{ROT}, etc.) are among
the fastest operations in StarForth --- they are direct array accesses on
the stack. Excessive stack manipulation is a sign of design issues, not a
performance problem per se. If you find yourself writing long sequences of
\texttt{ROT}, \texttt{SWAP}, and \texttt{OVER} to get values in the right
positions, consider using named variables instead:

\begin{lstlisting}[language=Forth]
\ Tangled stack version
: COMPLEX  ( a b c d -- result )
    ROT ROT OVER * ROT ROT ... ;

\ Cleaner with locals / variables
VARIABLE _A  VARIABLE _B  VARIABLE _C  VARIABLE _D

: COMPLEX  ( a b c d -- result )
    _D !  _C !  _B !  _A !
    _A @ _B @ * _C @ _D @ * + ;
\end{lstlisting}

The variable version is slightly slower per operation but far easier to read
and debug.

%%────────────────────────────────────────────────────────────────────────────
\chapter{The Block Editor}
\label{chap:block-editor}
%%────────────────────────────────────────────────────────────────────────────

%% ─────────────────────────────────────────────────────────────────────────
\section{Blocks as Source Storage}
\label{sec:blocks-as-source}

Traditional FORTH stored all source code in 1024-byte blocks on disk. StarForth
preserves this model. A block holds up to 16 lines of 64 characters each, no
newlines --- a fixed-format screen of text.

To list the contents of a block:

\begin{lstlisting}[language=Forth]
ok> 2048 LIST
\end{lstlisting}

This prints the 16 lines of block 2048, line-numbered. Empty blocks print as
blank lines.

To load (execute) a block:

\begin{lstlisting}[language=Forth]
ok> 2048 LOAD
\end{lstlisting}

To load a range of blocks:

\begin{lstlisting}[language=Forth]
ok> 2048 2055 THRU
\end{lstlisting}

%% ─────────────────────────────────────────────────────────────────────────
\section{Block Storage Backends}
\label{sec:block-backends}

StarForth supports three block storage backends:

\begin{description}
  \item[RAM] Default for interactive sessions. Blocks live in VM memory.
    Fast but not persistent across sessions.
  \item[File] Blocks are stored in a file on the host filesystem. Persistent.
    Configured by \texttt{STARFORTH\_BLOCK\_FILE} environment variable.
  \item[Factory] Read-only blocks baked into the binary. Used for capsule
    content in the bare-metal kernel.
\end{description}

On a standard hosted build, the default is RAM. To use file-backed blocks:

\begin{lstlisting}[language=bash]
STARFORTH_BLOCK_FILE=myblocks.blk ./build/amd64/standard/starforth
\end{lstlisting}

%% ─────────────────────────────────────────────────────────────────────────
\section{Protecting Blocks with FLUSH}
\label{sec:flush-protection}

If you have written to file-backed blocks and want to ensure changes are
saved before exiting:

\begin{lstlisting}[language=Forth]
ok> FLUSH
ok> BYE
\end{lstlisting}

\texttt{FLUSH} writes all modified (dirty) block buffers back to the backing
store. Exiting without \texttt{FLUSH} may lose changes if the backend does
not auto-flush on close.

%%────────────────────────────────────────────────────────────────────────────
\chapter{Capsules and the Boot System}
\label{chap:capsules}
%%────────────────────────────────────────────────────────────────────────────

%% ─────────────────────────────────────────────────────────────────────────
\section{What a Capsule Is}
\label{sec:capsule-what}

A capsule is a self-contained FORTH source file whose identity is its content
hash. The capsule ID is the XXHash64 of the capsule's contents --- any change
to the file produces a different ID. This makes capsules tamper-evident:
load a capsule, compute its hash, compare to the expected ID, and you know
whether the content is unchanged.

Capsule files live in the \texttt{capsules/} directory and use the
\texttt{.4th} extension.

%% ─────────────────────────────────────────────────────────────────────────
\section{init.4th: The Default Personality}
\label{sec:init-4th}

When StarForth starts, it loads \texttt{init.4th} from \texttt{./conf/} (on a
hosted build). This file defines the foundational word set, installs a
dictionary fence, and zeros blocks for user use.

The fence means that words defined in \texttt{init.4th} cannot be removed with
\texttt{FORGET}. This protects the core vocabulary from accidental deletion.

%% ─────────────────────────────────────────────────────────────────────────
\section{Capsule Variants}
\label{sec:capsule-variants}

StarForth ships several \texttt{init.4th} variants for different operating
modes:

\begin{center}
\begin{tabular}{ll}
\toprule
File & Purpose \\
\midrule
\texttt{init.4th}          & Default Mama VM personality \\
\texttt{init-0.4th} through \texttt{init-9.4th} & Numbered personality variants \\
\texttt{init-l8-stable.4th}   & L8 Jacquard stable mode \\
\texttt{init-l8-volatile.4th} & L8 Jacquard volatile mode \\
\texttt{init-l8-diverse.4th}  & L8 Jacquard diverse mode \\
\texttt{init-l8-temporal.4th} & L8 Jacquard temporal mode \\
\texttt{init-l8-transition.4th} & L8 Jacquard transition mode \\
\texttt{init-l8-omni.4th}    & L8 Jacquard omni mode \\
\texttt{ACL.4th}             & Word-level ACL subsystem (opt-in) \\
\texttt{zuse.4th}            & Bootstrap superuser (loaded by ACL.4th) \\
\bottomrule
\end{tabular}
\end{center}

The L8 Jacquard variants configure the Steady-State Machine (\texttt{ssm\_jacquard.c})
to operate in a specific mode. Each mode biases the adaptive runtime toward
different behavior: stable prefers low-variance execution, diverse favors
broad word coverage, temporal weights recent history more heavily, and so on.
These are intended for DoE experimentation rather than everyday use.

%% ─────────────────────────────────────────────────────────────────────────
\section{The Birth Protocol}
\label{sec:birth-protocol}

In the bare-metal kernel (LithosAnanke), StarForth VMs are born through a
birth protocol:

\begin{enumerate}
  \item Locate the capsule by name.
  \item Validate the content hash.
  \item Allocate a VM ID.
  \item Execute the capsule's IDENTITY block (the capsule code).
  \item Execute the PERSONALITY block (block 1 from the ramdrive).
  \item Log a parity record: VM ID + capsule hash + dictionary hash.
\end{enumerate}

The parity record enables offline determinism verification: an independent
system can load the same capsule and compare dictionary hashes. Identical
hashes mean identical dictionaries --- the definition of 0\% algorithmic
variance.

For hosted interactive use, the birth protocol runs transparently when
StarForth starts. You only interact with its results (the loaded dictionary),
not the protocol itself.

%%────────────────────────────────────────────────────────────────────────────
\chapter{The L8 Jacquard Mode Selector}
\label{chap:jacquard}
%%────────────────────────────────────────────────────────────────────────────

%% ─────────────────────────────────────────────────────────────────────────
\section{What the Jacquard Selector Does}
\label{sec:jacquard-what}

The L8 Jacquard Mode Selector is an additional layer above the seven feedback
loops. It monitors the distribution of execution patterns across attractor
buckets (clusters of similar execution histories) and switches between six
predefined mode configurations based on which attractor the system is
approaching.

Named after Joseph Marie Jacquard, whose programmable loom used punch cards to
control complex patterns, the Jacquard Selector treats the adaptive runtime
as a programmable pattern --- selecting different behavior profiles depending
on the observed workload dynamics.

%% ─────────────────────────────────────────────────────────────────────────
\section{The Six Modes}
\label{sec:jacquard-modes}

\begin{description}
  \item[\textbf{stable}] Favors low-variance, predictable execution. Decay
    slopes are conservative; the hot-words cache changes slowly. Best for
    long-running workloads with stable access patterns.
  \item[\textbf{volatile}] Responds quickly to changes in access patterns.
    Decay is aggressive; the cache refreshes rapidly. Best for workloads
    with rapidly shifting hot paths.
  \item[\textbf{diverse}] Biases toward broad dictionary coverage. Prevents
    a small set of words from monopolizing the cache at the expense of
    moderately-used words. Best for exploratory interactive sessions.
  \item[\textbf{temporal}] Weights recent execution history more heavily.
    The rolling window shrinks to focus on the most recent activity. Best
    for workloads with clear phases.
  \item[\textbf{transition}] Intermediate mode used during mode changes.
    Smooths the transition between attractor basins.
  \item[\textbf{omni}] All attractor statistics are weighted equally. No
    single mode dominates. Used when the workload does not fit a single
    characterization.
\end{description}

%% ─────────────────────────────────────────────────────────────────────────
\section{Mode Selection for Users}
\label{sec:jacquard-user}

In everyday interactive use you will not configure the Jacquard Selector
directly. It operates automatically based on workload statistics. The L8
mode variants of \texttt{init.4th} (\texttt{init-l8-stable.4th}, etc.)
pre-configure the Selector to start in a specific mode for DoE experiments
where you want to observe behavior in a controlled attractor.

If you are running a long benchmark and want to observe which mode the system
is in, use \texttt{WORD-ENTROPY} --- the output includes SSM state
information in builds where that diagnostic is enabled.

%%────────────────────────────────────────────────────────────────────────────
\chapter{Formal Verification}
\label{chap:formal-verification}
%%────────────────────────────────────────────────────────────────────────────

%% ─────────────────────────────────────────────────────────────────────────
\section{What Has Been Proved}
\label{sec:proofs-overview}

StarForth ships 24 Isabelle/HOL theory files providing machine-checkable
proofs of correctness and determinism. These proofs live in the \texttt{proof/}
directory and cover:

\begin{itemize}
  \item All seven feedback loops (Loop~\#1 through Loop~\#7), proving that
    each accumulation--stabilization pair behaves as designed.
  \item Core word categories: arithmetic, stack operations, logical operations,
    memory operations, return-stack operations, and Q48.16 fixed-point.
  \item Overall correctness (\texttt{StarForth\_Correctness.thy}), concurrency
    properties (\texttt{StarForth\_Concurrent.thy}), state transitions, and
    mutual exclusion.
  \item ACL system (five additional theories): pin monotonicity, inheritance
    clearing pin, TTL boundedness, emergency bypass, and no-escalation.
\end{itemize}

%% ─────────────────────────────────────────────────────────────────────────
\section{What Determinism Means Here}
\label{sec:determinism}

The proofs establish that the adaptive runtime is \emph{algorithmically
deterministic}: given the same sequence of word executions and the same
initial state, the runtime produces the same internal state every time.
This property was validated experimentally across 90 runs with 0\% variance.

Algorithmic determinism does not mean the same wall-clock times --- timing
depends on the host hardware, OS scheduler, and other factors outside StarForth's
control. It means the logic of adaptation (which words are promoted to the cache,
what decay slope is computed, what window size is selected) is fully determined
by the execution history, not by any random or non-deterministic process.

This matters for reproducible deployments: two instances of StarForth loading
the same capsule and receiving the same inputs will arrive at the same
dictionary heat distribution.

%% ─────────────────────────────────────────────────────────────────────────
\section{Checking the Proofs}
\label{sec:checking-proofs}

To check the Isabelle/HOL proofs (requires an Isabelle installation):

\begin{lstlisting}[language=bash]
isabelle build -D proof/
\end{lstlisting}

See \texttt{docs/01-getting-started/DEVELOPER.md} for Isabelle installation
instructions and compatible versions.

%%────────────────────────────────────────────────────────────────────────────
\chapter{StarForth on L4Re and Embedded Targets}
\label{chap:l4re}
%%────────────────────────────────────────────────────────────────────────────

%% ─────────────────────────────────────────────────────────────────────────
\section{L4Re Microkernel}
\label{sec:l4re}

StarForth can run on the L4Re microkernel operating system. The L4Re build
target enables the L4Re block-storage backend and ROMFS file access in place
of POSIX:

\begin{lstlisting}[language=bash]
make l4re
\end{lstlisting}

The L4Re build implies \texttt{STARFORTH\_MINIMAL=1}: ANSI color codes are
suppressed and I/O is reduced to minimal output. The build links without the
standard C library (\texttt{-nostdlib -ffreestanding}).

Status: partially implemented. The ROMFS stub is present; full L4Re integration
is ongoing.

%% ─────────────────────────────────────────────────────────────────────────
\section{Minimal / Freestanding Builds}
\label{sec:minimal-build}

The minimal build produces a freestanding binary with no libc dependencies:

\begin{lstlisting}[language=bash]
make minimal
\end{lstlisting}

This is useful for embedded environments or as a starting point for porting to
a new operating system. The minimal build disables features that depend on POSIX
or the C library. You can add features back selectively by enabling the
corresponding build flags.

%% ─────────────────────────────────────────────────────────────────────────
\section{Cross-Compilation}
\label{sec:cross-compile}

Cross-compile for AArch64 from an x86\_64 host:

\begin{lstlisting}[language=bash]
sudo apt-get install gcc-aarch64-linux-gnu
make CC=aarch64-linux-gnu-gcc rpi4-cross
scp build/starforth user@target-device:~/
\end{lstlisting}

The cross-compiled binary is statically linked and requires no runtime
dependencies on the target. Verify the architecture:

\begin{lstlisting}[language=bash]
file build/starforth
# ELF 64-bit LSB executable, ARM aarch64, statically linked
\end{lstlisting}

Cross-compilation for RISC-V 64 follows the same pattern with the appropriate
cross-toolchain (\texttt{gcc-riscv64-linux-gnu}).
